\section{成本下降、Jevons 悖论与产业走私}

\subsection{1200 倍的成本下降与 Jevons 悖论}

GPT-3 时代的推理成本是 \$60-70/百万 token。三年后（通过 Llama 3B 等小模型），同等能力降至几美分——\textbf{1200 倍}的降幅。DeepSeek 在 GPT-4 级别模型的成本趋势线上，但它是第一个到达这个价格点的。

NVIDIA 股价在 DeepSeek 发布后暴跌，但这是基于错误叙事的社交传染。没有任何公开模型的训练花费超过 \$1B。``DeepSeek 只花了 600 万美元''的说法是对论文的误读——那只是一次预训练运行的 GPU 小时成本，不包括研发、实验、后训练、推理部署。

\begin{importantbox}{Jevons 悖论在 AI 中的完美验证}
经济学中的 Jevons 悖论（1865 年提出）：当一种资源的使用效率提高时，总消耗量反而增加——因为降低成本刺激了更多需求。

蒸汽机效率提高 $\to$ 煤炭消耗不降反升。AI 推理成本降低 $\to$ GPU 需求不降反升。

DeepSeek 发布后的直接证据：
\begin{itemize}
    \item AWS H100 租用价格不降反\textbf{升}
    \item H200 几乎脱销
    \item 所有云厂商报告 GPU 需求激增
    \item 新的应用场景（推理密集型 Agent、大规模 RAG）变得经济可行
\end{itemize}

\textit{``'Scaling Laws 已死'持续了一个月……然后 O1、O3、R1 出来了，现在变成'模型进步太快了'。''} (Dylan)

更便宜的推理 $\to$ 更多的应用场景 $\to$ 更多的推理需求 $\to$ 需要更多 GPU。这不是零和博弈——总市场在扩大。这就是为什么 NVIDIA 的长期前景并不因为 DeepSeek 而黯淡。
\end{importantbox}

\subsection{GPU 走私的规模与手段}

GPU 走私已经从灰色地带发展为产业级操作，涉及多种渠道：

\textbf{云租用渠道}：字节跳动是最大的``走私者''——从 Oracle、Google、各小型云商全球租用 500,000+ GPU。通过分散在多个国家的实体来规避单一审查。这种``分布式走私''极难追踪，因为每个单独的租用合同都低于监管阈值。

\textbf{物理走私}：有人在旧金山机场头等舱带着 SuperMicro GPU 服务器箱子飞上海——这不是虚构，是真实发生的事件。2024 年估计有 200-300K GPU 通过新加坡、马来西亚的小公司流入中国。

\textbf{新扩散规则}的局限：限制中国关联实体租用 $<$2000 GPU 或购买 $<$1500 GPU。但执法挑战巨大——如何追踪全球云计算资源的最终用户？如何定义``中国关联实体''？一家在新加坡注册、由中国公民运营的公司算不算？

\textit{``GPU 将超越毒品和武器成为每公斤最高价值的走私品。''} (Dylan) 一张 H100 售价约 \$30,000，重量不到 5kg——每公斤价值 \$6,000+，远超大多数违禁品。而且不像毒品，GPU 是合法商品在某些市场，只是在特定出口方向上被限制。这创造了一个独特的套利空间。

\subsection{蒸馏、间谍与知识流动}

\subsubsection{蒸馏争议的本质}

蒸馏——用更强模型的输出来训练弱模型——是行业标准做法。OpenAI 声称 DeepSeek 使用了他们的 API 输出，金融时报大标题报道。但 Nathan 指出这更像是``叙事控制''而非真正的技术指控：

\textit{``为什么我用你的模型输出训练是不道德的，而你可以用互联网的文字训练？''} (Nathan) 用户通过 ChatGPT 生成的内容被复制粘贴到互联网上，成为了公开数据。很多美国创业公司公开在 OpenAI 输出上训练——为什么他们可以而 DeepSeek 不行？

更深层的问题是：\textbf{模型输出的知识产权归谁？}如果我用 GPT-4 生成了一段代码，这段代码的版权属于 OpenAI 还是属于我？这个问题在法律上尚无定论。

\subsubsection{人才流动与工业间谍}

硅谷的跳槽文化（加州非竞争条款违法）使得想法——不是代码——自由流动。一个具体案例：帮助 Google 建造 Gemini 100 万 context 的工程师跳槽到 Meta——下一代 Llama 预期将有 100 万 context。这是合法的知识转移还是事实上的技术盗窃？边界极其模糊。

代码/数据窃取是犯罪行为且相对容易追踪。但\textbf{想法的转移}是不可避免的——你不能``忘记''你在前公司学到的架构直觉。

\begin{knowledgebox}{日本：意外的 AI 训练天堂}
日本可能成为全球 AI 训练的一个独特节点：
\begin{itemize}
    \item \textbf{版权法}：明确允许在任何数据上训练模型（罕见的法律豁免）
    \item \textbf{能源}：9GW 搁置核电容量可重启用于数据中心——这是巨大的廉价电力来源
    \item \textbf{硬件}：无 GPU 进口限制（不受美国出口管制约束，因为日本是盟国）
    \item 多家日本公司和外资已开始在日本建设大型 AI 训练设施
\end{itemize}

版权豁免 + 廉价核电 + 无限 GPU = 理想的训练环境。这是一个很少人注意到的地缘套利机会。
\end{knowledgebox}

针对 AI 研究者的情报活动也在发生。Dylan 半开玩笑地指出：\textit{``作为一个二十多岁的单身男人……我们非常容易被腐蚀。''} 钓鱼行动（honey pot）针对拥有前沿知识的年轻研究者是真实存在的安全威胁。各国情报机构都在争夺 AI 人才——不是为了招募间谍，而是为了获取训练秘诀和架构洞察。

\subsection{Nathan 的 AI2/Tulu 开源使命}

Nathan 在 AI2（Allen Institute for AI）的工作代表了开源 AI 的最高标准。他的项目 Tulu 是第一个在 Chatbot Arena 排名前 60 的模型中公开\textbf{完整}后训练配方（代码 + 数据）的。

关键成果：Tulu 在 Llama 基座模型上应用 RLVR（可验证奖励的 RL），在数学推理上击败了 Llama 官方 instruct 版本，并匹配了 DeepSeek V3 的表现。这证明了\textbf{后训练方法本身}（而非仅仅是基座模型规模）可以带来巨大的性能差异。

OLMo 是 AI2 的另一个项目——完全开放的预训练工作（与 Tulu 的后训练互补）。更大模型规模下 RL 训练更容易激发强大能力，然后可以蒸馏到小模型。

Nathan 的核心动机：\textit{``我不信任那些说'相信我兄弟，我们会让 AI 变好'的人。''} AI 是我们一生中最强大的技术——需要更多人参与塑造它，而不是信任少数公司闭门决策。开放性帮助：非 AI 领域的研究者、政府、所有人都能理解正在发生什么。

\subsection{本章小结}

AI 产业的经济学正在极速重构：1200 倍的成本下降被 Jevons 悖论抵消，总需求持续增长。GPU 走私成为新的地下产业。蒸馏和人才流动使得知识产权边界极度模糊。日本的独特法律+能源+硬件组合创造了意外的训练优势。Nathan 的 AI2/Tulu 工作证明了完全开放的后训练可以达到前沿水平。

